\documentclass[11pt]{article}
\usepackage[margin=1in]{geometry}
\usepackage[T1]{fontenc}
\usepackage{lmodern,microtype}
\usepackage{amsmath,amssymb,amsthm,mathtools}
\usepackage{booktabs,longtable,array,enumitem}
\usepackage{listings}
\usepackage[colorlinks=true,linkcolor=blue,citecolor=blue,urlcolor=blue]{hyperref}
\hypersetup{pdftitle={Tape merge is optimal up to a size ratio of 148/95},pdfauthor={Yunjiang Jiang}}
\newtheorem{theorem}{Theorem}[section]
\newtheorem{lemma}[theorem]{Lemma}
\newtheorem{proposition}[theorem]{Proposition}
\newtheorem{corollary}[theorem]{Corollary}
\theoremstyle{definition}
\newtheorem{definition}[theorem]{Definition}
\newtheorem{example}[theorem]{Example}
\theoremstyle{remark}
\newtheorem{remark}[theorem]{Remark}
\newcommand{\D}{D}
\newcommand{\Z}{\mathbb Z}
\newcommand{\N}{\mathbb N}
\newcommand{\ceil}[1]{\left\lceil#1\right\rceil}
\newcommand{\floor}[1]{\left\lfloor#1\right\rfloor}
\newcommand{\rev}{\operatorname{rev}}
\newcommand{\sgn}{\operatorname{sgn}}
\lstset{basicstyle=\ttfamily\footnotesize,breaklines=true,columns=fullflexible,keepspaces=true,showstringspaces=false,frame=single}
\setlist{topsep=4pt,itemsep=3pt}
\title{Tape merge is optimal up to a size ratio of $148/95$}
\author{Yunjiang Jiang}
\date{3 October 2026}
\begin{document}
\maketitle
\begin{abstract}
Let $M(m,n)$ be the minimum worst-case number of pairwise comparisons needed
to merge two sorted lists of distinct elements, of lengths $m$ and $n$.
We give a computer-assisted proof that $M(m,n)=m+n-1$ whenever
$1\le m\le n\le(148/95)m$.
The constant $148/95=1.5578947368\ldots$ improves on the $38/25=1.52$
bound of Li, Sun and Zhang.
We retain Knuth's disjunctive adversary and extend the block-translation
argument of Li, Sun and Zhang. A parameterized induction reduces the
all-size statement to 138,622 integer inequalities, seven boundary values,
and 10,651 finite seed pairs. A two-direction block decomposition removes
the need for monotonicity lemmas in the final passage to all list sizes.
Two separately implemented recurrences agree on every entry of the required
finite table. We also give independently checked decision-tree certificates
for $M(6,22)=21$ and $M(6,25)=22$, values conjectured by Smith and Lang in 1994.
A secondary corollary of published upper bounds shows
$M(m,2m-3)\le3m-5$ for $m\ge9$.
All recurrences, exceptional cases, and verification procedures are specified.
\end{abstract}
\tableofcontents

\section{Introduction}
\label{sec:introduction}

\subsection{A simple algorithm and a difficult optimality question}
Suppose that
\[
 a_1<a_2<\cdots<a_m,\qquad b_1<b_2<\cdots<b_n
\]
are already sorted. All $m+n$ elements are distinct. Their values can be
accessed only through comparisons: a query asks whether $a_i<b_j$ or
$a_i>b_j$. Comparisons within either list give information already known,
so only comparisons between lists are useful. The task is to determine the
single sorted order of all the elements.

An algorithm may choose each query using the answers to earlier queries.
Its worst-case cost is the greatest number of queries it makes on any
possible interleaving of the two lists. We consider deterministic algorithms
and count comparisons, without imposing a bound on other computation or
memory. The smallest possible worst-case cost is denoted by $M(m,n)$.
By exchanging the lists, $M(m,n)=M(n,m)$; we usually take $m\le n$.
If one list is empty, no comparison is needed, so $M(m,0)=M(0,n)=0$.

Such an algorithm can be drawn as a binary decision tree: an internal
node specifies a comparison, its two branches are the possible answers,
and a leaf specifies the final merged order. An execution's depth is its
number of comparisons; the tree's height is the largest such depth.
We write $\floor{x}$ and $\ceil{x}$ for the largest integer at most $x$
and the smallest integer at least $x$, respectively. All logarithms have
base two.

The familiar \emph{tape merge} repeatedly compares the first remaining
element of each list and outputs the smaller one. Each comparison removes
one element. Once one list is empty the other can be appended, so
\begin{equation}\label{eq:tapeupper}
 M(m,n)\le m+n-1\qquad(m,n\ge1).
\end{equation}
Its worst case attains this bound. Optimality asks whether \emph{every}
correct algorithm has some input on which it needs this many comparisons.
For very unequal list lengths the answer is no: for $m=1$, binary search
in the $n+1$ possible insertion positions gives
$M(1,n)=\ceil{\log_2(n+1)}$.

There are exactly $\binom{m+n}{m}$ possible interleavings, because choosing
the $m$ positions occupied by the $a$-elements determines the entire order.
A binary decision tree of height $h$ has at most $2^h$ leaves; distinct
interleavings need distinct leaves. Consequently
\begin{equation}\label{eq:information}
 \ceil{\log_2\binom{m+n}{m}}\le M(m,n).
\end{equation}
This counting argument alone generally leaves a gap. Queries are restricted
to comparing two elements; they cannot ask an arbitrary yes/no question
that divides the remaining interleavings in half. The issue is therefore
both how many possibilities remain and how comparisons can separate them.

\subsection{Historical context and scope}
Graham and Karp independently established the optimality of tape merge for
equal lengths; Knuth records this result and develops the adversary method
used here~\cite[Section~5.3.2]{knuth}.
Hwang and Lin introduced binary merge, which adapts to unequal list
lengths~\cite{hwanglin1972}; they also determined $M(2,n)$~\cite{hwanglin1971}.
The cases of three and four elements in the smaller list required further
work by Hwang, Murphy, and M\"onting~\cite{hwang1980,murphy1978,monting1981}.
M\"onting also obtained substantial results for five elements.

The first broad tape-optimality region was
\[
 m\le n\le\floor{\frac32m}+1.
\]
It was proved independently by Stockmeyer and Yao~\cite{stockmeyeryao1980},
Murphy and Paull~\cite{murphypaull1979}, and Christen; the last result is
cited as an unpublished 1978 manuscript by Li, Sun and
Zhang~\cite{christenstraight,lsz2020}.
For lists farther from equal size, Christen's forward-testing and
backward-insertion method and subsequent improvements to binary merge
provide different upper bounds~\cite{christen1978,manacher1979,thanhbui1982}.
These algorithmic results are relevant background, but their procedures are
not used in the proof of our main theorem.

Smith and Lang's 1994 report~\cite{smithlang} used exact game search and
computable adversary recurrences to obtain many finite values and bounds.
One striking value is $M(7,12)=17$: tape merge would use $18$ comparisons.
Their work also makes clear why a large numerical table is not, by itself,
a theorem about all list sizes.
Li, Sun and Zhang~\cite{lsz2016,lsz2020} crossed the asymptotic ratio $3/2$
by proving tape optimality for $n\le(38/25)m$.
Their proof translates a finite adversary block of size $(25,38)$ into
arbitrarily large instances. They also prove that the same disjunctive
adversary cannot certify tape optimality when $n\ge9\ceil{m/5}$, and give
improved algorithms near $n=2m$.

Other merging models have different objectives. The randomized expected
comparison bounds of Fern\'andez de la Vega, Kannan and Santha~\cite{fernandez1993}
do not give deterministic worst-case bounds. Likewise, parallel depth,
in-place storage, and the total running time of a mergesort are different
questions from the function $M$ considered here. Linial's result on merging
under a known partial order~\cite{linial1984} is particularly relevant to
our finite exact computation, although that theorem is not needed by our
verifier.

\subsection{Results and attribution}
\begin{theorem}\label{thm:main}
For all positive integers $m,n$,
\[
 m\le n\le\frac{148}{95}m
 \quad\Longrightarrow\quad M(m,n)=m+n-1.
\]
\end{theorem}
The gain over $38/25$ is $18/475$ in the ratio.
To our knowledge, no stronger general ratio bound has been published;
the literature review for this manuscript was completed on 3 October 2026.
For example, at $m=1000$ the new bound reaches $n=1557$, whereas $38/25$
reaches $n=1520$. This is an extension of the Li--Sun--Zhang method within
Knuth's existing adversary framework. It does not remove the framework's
separate obstruction near ratio $1.8$.
The earlier additive region $n\le\floor{3m/2}+1$ may still be stronger for
some small $m$; the two results can be combined by taking their union.

\begin{theorem}\label{thm:exact}
\[M(6,22)=21,\qquad M(6,25)=22.\]
\end{theorem}
Smith and Lang already proved the corresponding upper bounds and explicitly
conjectured both equalities~\cite[Sections~3.12 and~4.3]{smithlang}.
Our certificates verify unrestricted lower bounds as well as matching
upper bounds. These assertions concern all legal comparison algorithms,
not only the bounded-width search classes considered in part of their
report.

\begin{proposition}\label{prop:uppermain}
For every integer $m\ge9$, $M(m,2m-3)\le3m-5$.
Consequently tape merge is suboptimal at every $n\ge2m-3$.
\end{proposition}
This is a short consequence of Smith--Lang finite algorithms and the
Li--Sun--Zhang recurrence and upper estimates. We include it as an
explicit corollary, not as a new merging procedure or a matching lower
bound. It supplies the upper-bound half of a family conjectured by Smith
and Lang~\cite[Section~5.1]{smithlang}.

\subsection{How the main proof works}
The proof has a finite part and an infinite part, connected by an explicit
induction. The following outline explains why both parts are needed.
\begin{enumerate}
\item \textbf{Make the adversary recursive.}
An adversary answers comparisons so as to delay the algorithm, while keeping
at least one input consistent with every answer. We restrict it to answers
that separate the remaining task into two smaller tasks. Some separations
share a single extreme element. Keeping track of which end elements are
privately constrained leads to nine recursively defined quantities
$D_{\lambda\rho}(m,n)$. Their unconstrained member satisfies
$D_{00}(m,n)\le M(m,n)$.
\item \textbf{Find a block that pays for all its elements.}
For $A=95$, $B=148$, and $S=A+B=243$, we prove translations of the form
\[
 D_{\lambda\rho}(m+A,n+B)
 \ge D_{\lambda\rho}(m,n)+S.
\]
This is stronger than simply putting two hard instances side by side:
that operation would ordinarily lose a comparison at each boundary.
We must instead answer \emph{every possible first comparison}, including
comparisons that cross the proposed block boundary.
\item \textbf{Repair the exceptional cases.}
Six boundary types suffice by symmetry. Two tiny states do not obey the
translation. Endpoint-deletion cuts and a uniform analysis of the comparison
$a_2:b_1$ repair every place where one of them appears as a child. Children containing elements
from only one list are handled separately. This gives an induction from a
finite square to all dimensions.
\item \textbf{Fill the whole ratio region.}
Exchanging the lists gives translations by both $(A,B)$ and $(B,A)$.
Any integer pair in the cone between these directions is a nonnegative
integer combination of the two blocks plus a positive residual pair with
both coordinates below $S$. Verifying the residual pairs completes the
proof. This last step uses no monotonicity theorem or conjecture.
\end{enumerate}

The computer checks only finite integer assertions. The all-size induction,
its exceptional cases, and the residual-pair argument are proved below.
The main finite table is computed twice, by differently organized
implementations of the same explicitly stated recurrence. The exact
certificates for Theorem~\ref{thm:exact} use a separate unrestricted
state space and a standalone verifier. The algorithms, accepted certificates,
and reproduction instructions accompany the source of this article.


% Standalone manuscript fragment; requires amsmath and amssymb.
% This file defines the recurrence and proves its lower-bound interpretation.
\section{A recursive adversary and its meaning}
\label{sec:game-definition}

All keys are distinct.  The input consists of two chains
$a_1<\cdots<a_m$ and $b_1<\cdots<b_n$.  A merging decision tree must be
correct for every interleaving of these chains.  Let $M(m,n)$ be the minimum
worst-case number of comparisons in such a tree.  Comparisons within a
chain carry no information and may be omitted.

\paragraph{Private restrictions are not public promises.}
We refer to the two list labels as colors: a block is monochromatic
if it contains keys from only one list, and bichromatic if it contains
keys from both. Write $0,L,R$ for Knuth's dot, backslash, and slash;
the symbols $\lambda$ and $\rho$ take values in this three-element set.  For $m,n>0$ a left
status $L$ restricts the adversary to completions with $a_1<b_1$, and a left
status $R$ restricts it to completions with $a_1>b_1$.  At the right, $L$
means $a_m<b_n$ and $R$ means $a_m>b_n$.  Status $0$ imposes no restriction.
These restrictions are obligations on the adversary; they are not
comparison outcomes given to the algorithm.  An algorithm can use an
endpoint inequality once the record of public comparison answers implies it, but
cannot introduce it as a premise merely because the adversary's private
state carries that status.

Here is a precise way to distinguish these two notions.  Let $\mathcal W$
be the set of all chain interleavings, and let $\mathcal W_{\lambda\rho}$
be its subset satisfying the two endpoint restrictions.  Define
\[
 H_{\lambda\rho}(m,n)
 =\min_{T\text{ correct on }\mathcal W}
       \max_{w\in\mathcal W_{\lambda\rho}}\operatorname{depth}_T(w).
\]
The minimum is still over globally correct trees, even though the
objective only uses the restricted subset.  A public-promise problem
would instead require correctness only on that subset, and is different.
For example, $H_{LL}(1,1)=1$, whereas the corresponding public-promise
problem needs no comparison.  Likewise, the private state $LR(2,1)$ has
only the permitted order $a_1<b_1<a_2$, but both adjacent cross-chain
relations must be established publicly, and its cost is two.

We define $D_{\lambda\rho}$ below by the recursive disjunctive rule.  The
proof establishes
\begin{equation}
 D_{\lambda\rho}(m,n)\le H_{\lambda\rho}(m,n)\le M(m,n),
 \qquad D_{00}(m,n)\le M(m,n).
 \label{eq:disjunctive-soundness}
\end{equation}
Thus no identification of $D$ with the unrestricted hidden-constraint
minimax problem is required.  In particular, the recurrence defines the
value of this recursive counterstrategy, not an assertion that every
possible private adversary is captured by it.

\paragraph{Feasibility and terminal conventions.}
For positive dimensions, the only impossible endpoint states are
\begin{equation}
 m=1,\ (\lambda,\rho)=(L,R),\qquad
 n=1,\ (\lambda,\rho)=(R,L).
 \label{eq:infeasible-game-states}
\end{equation}
Indeed, left $L$ and right $R$ require the first and last elements of the
merged order to be a-elements, and left $R$ and right $L$ require them
both to be b-elements.  These requirements are feasible exactly when the
corresponding chain has at least two elements; all remaining endpoint
combinations are feasible.  Assign impossible states value $-\infty$ and
use $-\infty+x=-\infty$.  They are never zero-cost alternatives.

A one-chain problem is terminal and costs zero.  Its nonexistent endpoint
comparisons are not retained as unresolved restrictions.  When a private
cut produces an a-only prefix, the parent's left status must be $0$ or
$L$; for a b-only prefix it must be $0$ or $R$.  An a-only suffix permits
right status $0$ or $R$; a b-only suffix permits $0$ or $L$.  Compatible
restrictions are then dropped.  This is the legal normalization of
monochromatic children.

For writing the six published formulas compactly, introduce a padded
value $\widehat D$ equal to $D$ at positive dimensions and satisfying
\[
 \widehat D_{\lambda\rho}(m,0)
 =\widehat D_{\lambda\rho}(0,n)=0
 \quad\text{for every }\lambda,\rho.
\]
These padded entries are used only with the endpoint exclusion tests
below.  They do not make arbitrary constrained zero-axis states legal
games.  In particular, the unambiguous unconstrained convention is
$D_{00}(m,0)=D_{00}(0,n)=0$.  A feasible state with $m,n>0$ is not terminal,
even when its private restrictions select a unique interleaving.

\paragraph{Six responses to a query.}
Fix a feasible parent $(m,n,\lambda,\rho)$ with $m,n>0$, and a first
comparison $a_i:b_j$.  Intervals such as $a_{u:v}$ denote the corresponding
consecutive subsequence, empty when $u>v$.  In the following table the
statuses beside a child are its left and right private statuses.  The
unprimed rows answer $a_i<b_j$; the primed rows answer $a_i>b_j$.
\begin{center}
\small
\begin{tabular}{c|c|c|c}
type & left child & right child & shared key\\ \hline
$A$ & $(a_{1:k},b_{1:\ell-1});\ (\lambda,0)$
    & $(a_{k+1:m},b_{\ell:n});\ (0,\rho)$ & none\\
$B$ & $(a_{1:k},b_{1:\ell});\ (\lambda,L)$
    & $(a_{k+1:m},b_{\ell:n});\ (R,\rho)$ & $b_\ell$\\
$C$ & $(a_{1:k},b_{1:\ell-1});\ (\lambda,R)$
    & $(a_{k:m},b_{\ell:n});\ (L,\rho)$ & $a_k$\\
$A'$ & $(a_{1:k-1},b_{1:\ell});\ (\lambda,0)$
     & $(a_{k:m},b_{\ell+1:n});\ (0,\rho)$ & none\\
$B'$ & $(a_{1:k-1},b_{1:\ell});\ (\lambda,L)$
     & $(a_{k:m},b_{\ell:n});\ (R,\rho)$ & $b_\ell$\\
$C'$ & $(a_{1:k},b_{1:\ell});\ (\lambda,R)$
     & $(a_{k:m},b_{\ell+1:n});\ (L,\rho)$ & $a_k$
\end{tabular}
\end{center}
Each child is reindexed starting at one.  All left-only keys are privately
placed before all right-only keys.  For an overlap, the repeated key is
privately the maximum of the left child and the minimum of the right
child.  The displayed inner statuses impose exactly these requirements:
an a-pivot has right status $R$ on the left and left status $L$ on the
right; a b-pivot has right status $L$ and left status $R$.

The complete parameter and endpoint restrictions are
\begin{center}
\begin{tabular}{c|c|c|c}
type & range of $k$ & range of $\ell$ & excluded endpoints\\ \hline
$A$  & $i\le k\le m$ & $1\le\ell\le j$
     & $\ell=1,\lambda=R$ or $k=m,\rho=R$\\
$B$  & $i\le k\le m$ & $1\le\ell<j$
     & $\ell=1,\lambda=R$ or $k=m,\rho=R$\\
$C$  & $i<k\le m$ & $1\le\ell\le j$
     & $\ell=1,\lambda=R$ or $k=m,\rho=R$\\
$A'$ & $1\le k\le i$ & $j\le\ell\le n$
     & $k=1,\lambda=L$ or $\ell=n,\rho=L$\\
$B'$ & $1\le k\le i$ & $j<\ell\le n$
     & $k=1,\lambda=L$ or $\ell=n,\rho=L$\\
$C'$ & $1\le k<i$ & $j\le\ell\le n$
     & $k=1,\lambda=L$ or $\ell=n,\rho=L$
\end{tabular}
\end{center}
In addition, any positive-dimensional child forbidden by
\eqref{eq:infeasible-game-states} invalidates the entire response.
The strict bounds in rows $B,C,B',C'$ are essential: the queried keys
must not belong together to a child, and neither may be the repeated
pivot.  They also prevent a child from having the parent's total size.
Every child of an admissible row has strictly fewer than $m+n$ keys,
counting the pivot once in each child that contains it.

The corresponding continuation scores, in the same order, are
\begin{align*}
 A(k,\ell)&=\widehat D_{\lambda0}(k,\ell-1)
             +\widehat D_{0\rho}(m-k,n+1-\ell),\\
 B(k,\ell)&=\widehat D_{\lambda L}(k,\ell)
             +\widehat D_{R\rho}(m-k,n+1-\ell),\\
 C(k,\ell)&=\widehat D_{\lambda R}(k,\ell-1)
             +\widehat D_{L\rho}(m+1-k,n+1-\ell),\\
 A'(k,\ell)&=\widehat D_{\lambda0}(k-1,\ell)
              +\widehat D_{0\rho}(m+1-k,n-\ell),\\
 B'(k,\ell)&=\widehat D_{\lambda L}(k-1,\ell)
              +\widehat D_{R\rho}(m+1-k,n+1-\ell),\\
 C'(k,\ell)&=\widehat D_{\lambda R}(k,\ell)
              +\widehat D_{L\rho}(m+1-k,n-\ell).
\end{align*}
These are Smith--Lang~\cite{smithlang}, equations (22)--(27), with
$0,L,R$ substituted for their symbols.  Let
$\mathcal R_{\lambda\rho}(m,n;i,j)$ be the legal rows and parameters
specified above.  The exact recurrence being computed is
\begin{equation}
 D_{\lambda\rho}(m,n)
 =1+\min_{1\le i\le m,\,1\le j\le n}
       \max_{r\in\mathcal R_{\lambda\rho}(m,n;i,j)}
              \operatorname{score}(r).
 \label{eq:nine-state-recurrence}
\end{equation}
For a feasible parent and any query, the response set is nonempty.
To see this, choose any feasible full interleaving and put a simple cut
between the queried keys.  The resulting consecutive prefix and suffix
obey the required endpoint tests.  Thus the maximum is at least zero.
Replacing the zeroed exclusion terms of Smith--Lang's equation (28) by
discarded terms, or adding an explicit zero to this maximum, does not
change the value.  This observation does not permit an impossible
positive-dimensional child to contribute zero to an otherwise valid sum.

\paragraph{Why the two children can be charged independently.}
Splits and endpoint statuses remain private.  Subsequent comparisons
whose two keys occur in one child are delegated to that child.  Such a
child is unique: the children are either disjoint or share only one key,
whereas a comparison has two distinct keys.  Every other cross-chain
comparison is between a left-only and a right-only key and receives the
fixed left-before-right answer.  It is an additional paid comparison,
not a free disclosure of the private split.

There are two facts to prove.  First, independent child answers are
jointly consistent.  For a simple split, concatenate endpoint-compatible
child completions.  For an overlap, the pivot is last in the left
completion and first in the right completion; concatenate and identify
the two copies.  This gives an interleaving of the original chains,
satisfies the parent's endpoint statuses, and has the advertised answer
to the first query.  The pivot is one physical key, not a duplicated
independent input.

Second, fixed cross-child answers cannot replace an unresolved internal
comparison certificate.  Consider any alternative completion of one
child consistent with that child's public local transcript.  This
alternative need not obey its private endpoint statuses.  Keep a
completion of the sibling that does obey its statuses.  For a simple
split, concatenate the two orders.  For a pivot split, if the altered
child is the left one, concatenate its entire alternative order with the
right order after deleting the right order's first element, the pivot.
If the altered child is the right one, concatenate the left order after
deleting its last element, the pivot, with the entire alternative right
order.  In both constructions the original a-chain and b-chain remain
in order, the sibling's order is preserved, and every left-only key is
still before every right-only key.  All delegated and cross-child
comparison answers therefore remain true.  Comparisons involving the
pivot and another child key were delegated to that child and are already
included in its local transcript.

Consequently, if one child's local transcript allows two different
interleavings, the full transcript also allows two different
interleavings.  The altered full interleaving may violate a private
restriction, which is immaterial: the merging algorithm is required to
be correct on that interleaving too.  This replacement argument composes
up an arbitrary recursive split tree.  At each ancestor the sibling can
be kept at an endpoint-compatible completion, so all previously answered
cross comparisons are preserved.  Revealing the split together with its
order inequalities as a public promise would destroy this argument and
would define the wrong game.

\paragraph{Proof of the lower-bound interpretation.}
Induct on $m+n$.  The zero-axis cases are immediate.  At a feasible
positive state the public transcript initially contains only the two
chain orders, so a globally correct algorithm must make a cross-chain
comparison.  For its first query, the adversary chooses a legal response
maximizing the score in \eqref{eq:nine-state-recurrence} and then runs the
two inductively supplied child strategies.  All actual completions obey
the private statuses by the consistency argument.  The replacement
argument forces the algorithm to finish a complete public certificate
in each child before it can terminate.  The induction hypothesis charges
at least the first child's $D$-value to comparisons delegated there and
at least the second child's $D$-value to comparisons delegated there.
No comparison is charged twice, even for an overlapping split.  The first
query belongs to neither child and costs one additional comparison.
Any later cross-child comparisons only increase the total.
The chosen response therefore forces its score plus one.  Taking the
minimum over possible first queries proves
$D_{\lambda\rho}\le H_{\lambda\rho}$ for every feasible state, and
\eqref{eq:disjunctive-soundness} follows.  Tape merge supplies
$M(m,n)\le m+n-1$ for positive $m,n$.

\paragraph{Three geometric split forms and the primary implementation.}
The six rows can equivalently be organized into three forms, each with
two possible query rectangles:
\begin{align*}
\text{simple: }& (p,q),(m-p,n-q),\quad
     i\le p,\ j\ge q+1\quad\text{or}\quad i\ge p+1,\ j\le q;\\
\text{a-pivot: }& (p,q),(m-p+1,n-q),\quad
     i\le p-1,\ j\ge q+1\quad\text{or}\quad i\ge p+1,\ j\le q;\\
\text{b-pivot: }& (p,q),(m-p,n-q+1),\quad
     i\le p,\ j\ge q+1\quad\text{or}\quad i\ge p+1,\ j\le q-1.
\end{align*}
The simple form uses $0\le p\le m$, $0\le q\le n$ and excludes an empty
or full prefix.  Its statuses are $(\lambda,0),(0,\rho)$, normalized on
monochromatic children as above.  The genuine a-pivot form uses
$1\le p\le m$, $1\le q<n$ and statuses
$(\lambda,R),(L,\rho)$.  The genuine b-pivot form uses
$1\le p<m$, $1\le q\le n$ and statuses
$(\lambda,L),(R,\rho)$.

The substitutions linking the two organizations are
\begin{align*}
 A &: (p,q)=(k,\ell-1),& A' &: (p,q)=(k-1,\ell),\\
 B &: (p,q)=(k,\ell),& B' &: (p,q)=(k-1,\ell),\\
 C &: (p,q)=(k,\ell-1),& C' &: (p,q)=(k,\ell).
\end{align*}
The published padded formulas also include four degenerate pivot cases.
They are sound but unnecessary.  Dropping a private endpoint restriction
cannot decrease $D$: this follows directly by induction in the recurrence,
since every former response is still available and its affected child
value can only increase.  Therefore each degenerate case is dominated
by the following simple cut:
\begin{center}
\begin{tabular}{c|c|c}
degenerate row & monochromatic child & dominating simple cut $(p,q)$\\ \hline
$B$, $k=m$ & right b-only & $(m,\ell)$\\
$C$, $\ell=1$ & left a-only & $(k-1,0)$\\
$B'$, $k=1$ & left b-only & $(0,\ell-1)$\\
$C'$, $\ell=n$ & right a-only & $(k,n)$
\end{tabular}
\end{center}
In each row the nonzero child retains its dimensions and loses one
private endpoint restriction.  The table's strict query bounds make the
dominating simple cut applicable, and the endpoint exclusions make it
compatible with the parent.  This proves the equivalence between the
literal six-formula program and the primary three-form program, including
their different zero-axis representations.

\paragraph{Symmetries.}
Define $\bar0=0$, $\bar L=R$, $\bar R=L$.  Swapping the two lists, and
reversing the entire order while reversing each list, respectively give
\[
 D_{\lambda\rho}(m,n)=D_{\bar\lambda\bar\rho}(n,m),\qquad
 D_{\lambda\rho}(m,n)=D_{\bar\rho\bar\lambda}(m,n).
\]
Both maps preserve feasibility, normalized terminal cases, query
rectangles, and the complete set of legal response trees, so these are
exact equalities for the recurrence.  In particular,
$D_{0L}=D_{R0}$, $D_{0R}=D_{L0}$ and $D_{LL}=D_{RR}$ at fixed dimensions.
The reversal-fixed states are $00,LR,RL$; a convenient set of six
representatives is $00,L0,R0,LL,LR,RL$.  The list-swap symmetry additionally
allows the computation to be restricted to $m\le n$.



\section{A self-contained finite-to-infinite translation argument}
\label{sec:expanded-translation}

The purpose of this section is to explain why a finite exact computation can
prove an assertion for lists of every length. The computation checks a finite
collection of inequalities. The induction below then propagates them to all
sizes. The argument uses the Knuth disjunctive recurrence developed in
Smith--Lang and Li--Sun--Zhang \cite{smithlang,lsz2016,lsz2020}. No external
increment or monotonicity theorem is needed for the propagation argument.

We use the states, feasibility rules, and three geometric split forms
defined in Section~\ref{sec:game-definition}. Every child is nonempty and
strictly smaller than its parent: a simple cut divides $m+n$ elements into
two nonempty parts, while a genuine pivot split has two parts of size at
least two whose sizes sum to $m+n+1$. Hence neither pivot child has size
more than $m+n-1$. The six boundary representatives are
$00,R0,L0,RL,LL,LR$; reversal supplies the other three states.

\subsection{Four auxiliary facts, with proofs}

\paragraph{Constraint removal.}
Removing an endpoint condition can only increase $D$. Induct on total size.
Every split compatible with the stronger condition is compatible with the
weaker one. Its corresponding child has either the same conditions or a
weaker outer condition, so its value does not decrease by induction. The
maximum for each fixed query does not decrease, and neither does the
minimum of these query values. Terminal and infeasible states agree with
this argument under the conventions above.

\paragraph{Direct sums within the restricted recurrence.}
Take two nonempty blocks, of sizes $(p,q)$ and $(x,y)$, and require every
first-block element to be smaller than every second-block element.
Let $\lambda$ be a compatible outer left condition and $\rho$ a compatible
outer right condition. On the first block retain $\lambda$ if it is
bichromatic, or drop it if it is monochromatic and automatically satisfied;
on the second do the same with $\rho$. Denote the resulting block values
by $U,V$. Then
\begin{equation}
 D_{\lambda\rho}(p+x,q+y)\ge U+V.
 \tag{DS}
\end{equation}
For two bichromatic blocks this reads
$D_{\lambda\rho}(p+x,q+y)\ge
D_{\lambda0}(p,q)+D_{0\rho}(x,y)$.

We prove (DS) by induction on the sum of the two block sizes. A first query
with one element in each block has a legal simple response at their boundary,
with value $1+U+V$. If both queried elements lie in the first block, that
block is necessarily bichromatic. Use an optimal response in that block.
Keep its left child and append the entire second block to its right child.
The same split is admissible in the full problem: for a simple split its
cut coordinates $(p',q')$ are unchanged, and for a complex split its pivot
and cut coordinates are unchanged. The comparison inequalities listed above
are therefore unchanged. Its right child has the old inner left condition
and the full outer right condition $\rho$. Apply (DS) inductively to this
enlarged right child. This is legitimate even if the old right child is
monochromatic: its inherited inner condition is then either absent (simple
split) or, for a complex split, both colors are present. The second block
carries the compatible outer right condition, exactly as in the statement.
The value of the full reply is at least the optimal old reply value plus
$V$, and hence at least $U+V$.

If the first query is inside the second block, prepend the first block to
its left child. In that case shift the local cut and, if present, its pivot
coordinates by $(p,q)$. The local query indices shift by the same amounts,
so each admissibility inequality is preserved. The inherited outer left
condition is compatible with the first block by assumption. Again apply
induction to the enlarged child. In either construction the local child
has fewer elements than its local parent, so the two-block total decreases.
A monochromatic or singleton block has no internal cross-list query; thus
it causes no omitted base case. If the full problem is monochromatic, both
block values are zero and the claim is immediate. This proves (DS) purely
inside the restricted recurrence; no unrestricted direct-sum theorem is
being substituted for it.

Appending a monochromatic tail in (DS) gives, for $u,v\ge1$,
\begin{equation}
 D_{R0}(A,B+v)\ge D_{R0}(A,B+1),\qquad
 D_{L0}(A+u,B)\ge D_{L0}(A+1,B).
 \tag{AX}
\end{equation}
For zero-length tails the inequalities are equalities. Prepending a
monochromatic block similarly shows that $D_{0\rho}(x,y)$ cannot decrease
when either dimension is increased, provided the indicated states and
outer condition are compatible. Only this explicitly proved one-sided
monotonicity will be needed below.

\paragraph{Endpoint deletion.}
For feasible positive-dimensional states, querying $(1,1)$ gives
\begin{equation}
 D_{R\rho}(m,n)\le1+D_{0\rho}(m,n-1),\qquad
 D_{L\rho}(m,n)\le1+D_{0\rho}(m-1,n).
 \tag{ED}
\end{equation}
To verify the first bound, the left $R$ condition permits only the answer
$a_1>b_1$. The split inequalities rule out both complex types. A simple
response must have $p=0$, $q\ge1$, with a b-only prefix of value zero and
suffix value $D_{0\rho}(m,n-q)$. This suffix is bounded above by
$D_{0\rho}(m,n-1)$ by the preceding consequence of (DS). The second bound
is symmetric. A zero-dimensional suffix is evaluated by the explicit
compatibility convention, so (ED) does not assert an undefined comparison.
Reversing both lists gives the corresponding right-endpoint deletion bounds.

\paragraph{A two-element query bound.}
For $m\ge3$, $n\ge2$ and $\rho\in\{0,L,R\}$, put
\[
 X_\rho=D_{0\rho}(m-2,n),\qquad
 Y_\rho=D_{0\rho}(m-1,n-1).
\]
Whenever the indicated original state is feasible,
\begin{equation}
 D_{L\rho}(m,n)\le\max\{1+X_\rho,\ 2+Y_\rho\}.
 \tag{Q}
\end{equation}
This follows by evaluating every response to query $(2,1)$, not by
assuming that the adversary chooses a particular response.
For a simple response $a_2<b_1$, the cut has $p\ge2,q=0$. Its prefix is
a-only and has value zero. Its suffix, after increasing its a-dimension
if necessary, has value at most $X_\rho$. For a simple response $a_2>b_1$,
compatibility with left $L$ rules out $p=0$, so $p=1,q\ge1$. The prefix has
value $D_{L0}(1,q)=1$, and the suffix has value at most $Y_\rho$.
To see $D_{L0}(1,q)=D_{LL}(1,q)=1$, query $(1,1)$: only the simple cut with
a-only prefix and b-only suffix is admissible, and one comparison is
necessary in a bichromatic nonterminal state.

For a shared a-pivot, the answer $a_2<b_1$ is impossible since it would
require $1>q\ge1$. The other answer requires pivot $p=1$, whose prefix
$LR(1,q)$ is infeasible. For a shared b-pivot, the first answer is again
impossible. The other forces $p=1$ and $q\ge2$. Its prefix is
$LL(1,q)$ of value one. After dropping the left $R$ constraint of its
suffix, the suffix is bounded by
$D_{0\rho}(m-1,n-q+1)\le Y_\rho$. All dimension comparisons here are the
one-sided consequences of (DS), with compatible right conditions.
These cases exhaust both answers and all three split types, proving (Q).

\subsection{The seven boundary checks and the finite translation square}
Fix positive integers $A,B$, put $S=A+B$, and choose
\[
 H\ge\max\{A+4,B+3\}.
\]
Require these seven inequalities:
\begin{align*}
 D_{00}(A,B)&\ge S-1,& D_{R0}(A,B)&\ge S-1,\\
 D_{R0}(A,B+1)&\ge S,& D_{LR}(A+1,B)&\ge S-1,\\
 D_{L0}(A+1,B)&\ge S,& D_{L0}(A+2,B)&\ge S+1,\\
 D_{L0}(A+1,B+1)&\ge S+1.&&
 \tag{B}
\end{align*}
In addition, verify throughout $0\le m,n\le H$ that
\begin{equation}
 D_{\lambda\rho}(m+A,n+B)\ge D_{\lambda\rho}(m,n)+S
 \tag{T}
\end{equation}
for the following domains:
\begin{itemize}
\item $00$: $m,n\ge0$, $m+n\ge1$;
\item $R0,L0,RL$: positive dimensions and a feasible state;
\item $LL$: positive dimensions, excluding $(1,1)$;
\item $LR$: positive dimensions and a feasible state, excluding $(2,1)$.
\end{itemize}
The excluded whole states are genuine exceptions to the proposed
translation and are not silently included in the theorem.
The $00$ axes have value zero and must be checked if a stored table only
lists positive dimensions. The double-zero state is not a translation base.

\begin{theorem}[Finite translation criterion]
\label{thm:translation}
If (B) and the finite square of (T) hold, then (T) holds globally on exactly
these domains. This is a simultaneous statement for the six boundary types.
\end{theorem}

\subsection{Copying a split and locating every exceptional child}
Suppose a query in the old $(m,n)$ state has an optimal response with two
children. Appending $(A,B)$ to the right child keeps the old cut parameters
and any pivot unchanged. All query inequalities remain true. The old right
child gains $A$ a-elements and $B$ b-elements. Prepending to the left child
instead shifts cut parameters and any pivot from $(p,q)$ to $(p+A,q+B)$;
the full query is $(i+A,j+B)$, so the inequalities again remain true.
The untouched child has exactly its old dimensions and conditions.

A simple monochromatic child needs special treatment because its original
inherited condition was automatically dropped. A b-only prefix with inherited
left $R$ enlarges to $D_{R0}(A,B+v)\ge S$ by (AX) and (B); an a-only prefix
with left $L$ enlarges to $D_{L0}(A+u,B)\ge S$. The opposite colors would have
been incompatible with the original split. Reversal gives the corresponding
repairs for a b-only suffix with right $L$ and an a-only suffix with right $R$.
An unconstrained monochromatic child uses the $00$ axis translation.
All complex children are bichromatic. Therefore no copying step below uses
a constrained translation at an undefined axis state.

The following list derives directly from the split rules every exceptional
child that can arise in the copy directions used in the six cases below.
Here $(i,j)$ are the \emph{old}, unshifted query indices. Other boundary
orientations are treated by reversal rather than by additional copy cases.
\begin{enumerate}
\item A prefix $RR(1,1)$ can only be an a-pivot split with $p=q=1$.
The answer $a_i<b_j$ would require $i<1$ and is impossible. The other
answer requires $i\ge2,j=1$. The pivot ranges and these query indices
force $m,n\ge2$.
\item A prefix $LL(1,1)$ can only be a b-pivot split with $p=q=1$.
The first answer gives $i=1,j\ge2$; the other would require $j<1$.
Here again the pivot ranges and query indices force $m,n\ge2$.
\item A prefix $LR(2,1)$ can only be an a-pivot split with $p=2,q=1$.
The two answer regions are $i=1,j\ge2$ and $i\ge3,j=1$.
Since $q<n$, this exception requires $n\ge2$.
\item A suffix $LL(1,1)$ can only be an a-pivot split with $p=m,q=n-1$.
Its first answer requires $i\le m-1,j=n$ and its other answer would
require $i>m$. In particular $m,n\ge2$.
\end{enumerate}
A simple split has an inner $0$ condition and thus cannot create any of
these doubly constrained exceptions. The list is exhaustive.

We shall repeatedly use this uniform repair for an old state $L\rho$,
where $\rho\in\{0,L,R\}$, when the full query satisfies
$i\ge A+3,j=B+1$ and has arisen from a prefix $LR(2,1)$.
That geometry ensures $m\ge3,n\ge2$, so (Q) applies. Set
$X=X_\rho$, $Y=Y_\rho$. If $X\ge1+Y$, (DS) at the block sizes
$(A+2,B)$ and $(m-2,n)$ yields
\begin{equation}
 D_{L\rho}(m+A,n+B)
 \ge D_{L0}(A+2,B)+X\ge S+1+X\ge S+D_{L\rho}(m,n).
 \tag{Repair-1}
\end{equation}
This bounds the whole state, so it bounds the adversary value at every
first query. It does not presuppose a publicly revealed split or an extra
comparison at that boundary. If $X<1+Y$, use the simple cut $(A+1,B+1)$.
It is admissible for this full query because $i>A+1$ and $j\le B+1$.
Its value is
\begin{equation}
 1+D_{L0}(A+1,B+1)+Y\ge S+2+Y\ge S+D_{L\rho}(m,n).
 \tag{Repair-2}
\end{equation}
The same repair therefore covers each of the three possible right
conditions, including the absence of a right condition.

\subsection{Proof of the six global translations}
Induct simultaneously on $m+n$, using the finite square as the base.
Outside it, $m>H$ or $n>H$. A copied old child is always strictly smaller
than the old parent, as proved above. The induction applies to that child,
with the explicit monochromatic repairs if appropriate, unless it is one
of the exceptional types just listed. Thus neither a nonshrinking recurrence
call nor an excluded induction hypothesis is being assumed.

Write $D^{i,j}_{\lambda\rho}$ for the maximum response value at a fixed
first query. Showing a lower bound for every such value proves it for $D$.
A lower bound obtained directly for the whole state, as in (Repair-1), also
bounds each query value because $D$ is their minimum.

Two elementary coverage observations will be used repeatedly. The four
regions relative to the old corner $(m,n)$ are the old rectangle
$i\le m,j\le n$, the two crossing rectangles, and
$i\ge m+1,j\ge n+1$. In the last region, $i\le A,j\le B$ cannot both hold:
they would imply $m\le A-1,n\le B-1$, inside the finite base. Also, whenever
$i\ge A+1,j\ge B+1$, the indices $(i-A,j-B)$ belong to the old problem,
because $i\le m+A,j\le n+B$. Thus all copied old queries below exist.
Reversal in the enlarged problem always means the explicit transformation
\[
 i'=m+A+1-i,\qquad j'=n+B+1-j.
\]

\paragraph{Case (a): $00$.}
If $i\le m,j\ge n+1$, or $i\ge m+1,j\le n$, a simple split into $(m,n)$ and
$(A,B)$ has value at least
\[
 1+D_{00}(m,n)+D_{00}(A,B)\ge S+D_{00}(m,n).
\]
If $i\le m,j\le n$, copy an optimal old response and append $(A,B)$ to its
right child. Its right constraint is $0$, so the nonexceptional $00,L0,R0$
translations suffice. If $i\ge m+1,j\ge n+1$, the inequalities
$i\le A,j\le B$ cannot both hold outside the base square. If the reversed
query also had $i'\ge m+1,j'\ge n+1$, the displayed transformation would
give precisely $i\le A,j\le B$, a contradiction. Hence the reversed query
lies in one of the first three regions, already treated. State $00$ is
reversal-invariant. If the old state is on an axis, the old-rectangle query
case is empty, and the same crossing or reversal argument still applies;
the old block is nonempty because $(m,n)\ne(0,0)$.

\paragraph{Case (b): $R0$.}
If $i\le m,j\ge n+1$, or $i\ge m+1,j\le n$, split into the old block and
the added block. The value is at least
$1+D_{R0}(m,n)+D_{00}(A,B)\ge S+D_{R0}(m,n)$.
If $i\le m,j\le n$, append the added block to the right child of an optimal
old response. That child's right condition is $0$, so its translation is
one of $00,L0,R0$, including the stated axis repairs. No exceptional
double-boundary state arises. In the remaining upper-right region,
if $i\le A,j\ge B+1$ or $i\ge A+1,j\le B$, the simple cut gives
\[
 1+D_{R0}(A,B)+D_{00}(m,n)\ge S+D_{R0}(m,n).
\]
If $i\ge A+1,j\ge B+1$, copy the optimal old response to
$(i-A,j-B)$ and prepend the block to its left child. The only exception is
$RR(1,1)$, equivalent by reversal to $LL(1,1)$. Its geometry forces
$i\ge A+2,j=B+1$. Use instead
\[
 1+D_{R0}(A,B+1)+D_{00}(m,n-1)
 \ge S+1+D_{00}(m,n-1)\ge S+D_{R0}(m,n).
\]

\paragraph{Case (c): $L0$.}
For the two crossing regions at the old corner, the simple split has value
$1+D_{L0}(m,n)+D_{00}(A,B)\ge S+D_{L0}(m,n)$.
For $i\le m,j\le n$, append to the right child of an optimal old response;
again only $00,L0,R0$ translations and the axis repairs are needed.
In the remaining upper-right region, if $i\le A,j\ge B+1$ or
$i\ge A+2,j\le B$, repeat the pivot
$a_{A+1}$ to obtain
\[
 1+D_{LR}(A+1,B)+D_{L0}(m,n)\ge S+D_{L0}(m,n).
\]
The missing case $i=A+1,j\le B$ forces $m\le A,n\le B-1$ and is in the base.
For $i\ge A+1,j\ge B+1$, prepend to the left child of an optimal old response.
An exceptional child $LL(1,1)$ forces $i=A+1,j\ge B+2$, and is repaired by
\[
 1+D_{L0}(A+1,B)+D_{00}(m-1,n)\ge S+D_{L0}(m,n).
\]
This repair works for every query at that location, including a response
whose exceptional child is $LR(2,1)$.

The only remaining $LR(2,1)$ location is $i\ge A+3,j=B+1$.
Apply (Repair-1) or (Repair-2) with $\rho=0$. This includes $i=A+3$;
there is no omitted equality case at this endpoint.

\paragraph{Case (d): $RL$.}
If $i\le A,j\ge B+1$ or $i\ge A+1,j\le B$, use the simple cut
\[
 1+D_{R0}(A,B)+D_{0L}(m,n)\ge S+D_{RL}(m,n).
\]
For $i\ge A+1,j\ge B+1$, prepend to the old left child. Its only exception is
$RR(1,1)$, forcing $i\ge A+2,j=B+1$; replace it by
\[
 1+D_{R0}(A,B+1)+D_{0L}(m,n-1)\ge S+D_{RL}(m,n).
\]
For $i\le A,j\le B$, reverse both lists. The state $RL$ is reversal-invariant,
and $i'=m+A+1-i,j'=n+B+1-j$. If also $i'\le A,j'\le B$, these equations
would imply $m+1\le i\le A,n+1\le j\le B$, forcing
$m\le A-1,n\le B-1$, inside the base. Thus the reversed query belongs
to one of the other three regions at the $(A,B)$ corner, already covered.

\paragraph{Case (e): $LL$.}
The regions $i\le m,j\ge n+1$ and $i\ge m+1,j\le n$ admit the simple cut
\[
 1+D_{L0}(m,n)+D_{0L}(A,B)\ge S+D_{LL}(m,n).
\]
For $i\le m,j\le n$, append to the right child of an optimal old response.
Its only exception is $LL(1,1)$; the query then has $i\le m-1,j=n$.
The repair is
\[
 1+D_{L0}(m,n-1)+D_{0L}(A,B+1)\ge S+D_{LL}(m,n).
\]
The last inequality uses right-endpoint deletion and retains the left $L$
constraint, as required by the child state.

In the upper-right case, the repeated-pivot split used in (c) gives
\[
 1+D_{LR}(A+1,B)+D_{LL}(m,n)\ge S+D_{LL}(m,n)
\]
when $i\le A,j>B$ or $i\ge A+2,j\le B$. The remaining case $i=A+1,j\le B$
is in the finite base. For $i\ge A+1,j\ge B+1$, prepend to the old left child.
The $LL(1,1)$ location $i=A+1,j\ge B+2$ is repaired by
\[
 1+D_{L0}(A+1,B)+D_{0L}(m-1,n)\ge S+D_{LL}(m,n).
\]
The same repair covers an $LR(2,1)$ child at that location.

The sole remaining $LR(2,1)$ location is $i\ge A+3,j=B+1$.
Use (Repair-1) or (Repair-2) with $\rho=L$. All their hypotheses follow
from the exceptional-child geometry already established.

\paragraph{Case (f): $LR$.}
This state is reversal-invariant. If $n>H$, choose the orientation with
\[
 j\ge\left\lceil\frac{B+n+1}{2}\right\rceil\ge B+2.
\]
For $i\le A$, repeat $a_{A+1}$ and obtain
$1+D_{LR}(A+1,B)+D_{LR}(m,n)\ge S+D_{LR}(m,n)$.
For $i\ge A+1$, prepend to the optimal old left child. The $LL(1,1)$ exception
forces $i=A+1,j\ge B+2$ and is repaired by
\[
 1+D_{L0}(A+1,B)+D_{0R}(m-1,n)\ge S+D_{LR}(m,n).
\]
That same cut covers any $LR(2,1)$ exception at this location. Its only
remaining location has $j=B+1$, excluded by the orientation.

If $n\le H$, then $m>H$. Choose the reversal with
\[
 i\ge\left\lceil\frac{A+m+1}{2}\right\rceil\ge A+2.
\]
For $j\le B$, repeat $a_{A+1}$ as above. For $j\ge B+1$, prepend to the old
left child. The $LL(1,1)$ exception would require $i=A+1$ and is impossible.
For the $LR(2,1)$ exception, the location $i=A+1$ is excluded by
orientation, so the only possibility is $i\ge A+3,j=B+1$.
Use (Repair-1) or (Repair-2) with $\rho=R$.

This completes the simultaneous induction. Every copied nonexceptional
child is strictly smaller, every monochromatic child was handled by an
axis repair, and every occurrence of either exceptional state was explicitly
repaired. The exceptional whole states are never invoked by induction.

\subsection{A finite cone of residual sizes}
The last step must not assume that tape optimality is monotone in one list
length. That unrestricted monotonicity is related to a classical open
conjecture. Instead we use the translation in both coordinate directions.
The next elementary lemma reduces every pair in a cone to a bounded pair
in the same cone.

\begin{lemma}[Finite cone remainder]
\label{lem:cone-remainder}
Let $1\le A<B$ and $S=A+B$. If positive integers $m,n$ satisfy
\[
 A m\le B n,\qquad A n\le B m,
\]
then there are nonnegative integers $k,\ell$ and positive integers $r,t<S$
such that
\[
 (m,n)=k(A,B)+\ell(B,A)+(r,t),
 \qquad A r\le B t,\quad A t\le B r.
\]
\end{lemma}
\begin{proof}
Put $\Delta=B^2-A^2>0$ and define
\[
 x=\frac{Bn-Am}{\Delta},\qquad
 y=\frac{Bm-An}{\Delta}.
\]
The cone inequalities imply $x,y\ge0$, and direct multiplication gives
$(m,n)=x(A,B)+y(B,A)$. Initially take $k=\lfloor x\rfloor$,
$\ell=\lfloor y\rfloor$, and set
\[
 (r,t)=(m,n)-k(A,B)-\ell(B,A)
      =\{x\}(A,B)+\{y\}(B,A),
\]
where $\{z\}=z-\lfloor z\rfloor$ is the fractional part. Thus $r,t$ are
integers in $[0,S)$ and are nonnegative combinations of the two bounding
vectors. They satisfy the same cone inequalities. If either fractional
part is nonzero, positivity of all four vector coordinates implies
$r,t>0$. If both fractional parts vanish, then $(r,t)=(0,0)$ and
$k+\ell>0$, since $m,n>0$. If $k>0$, replace $k$ by $k-1$ and $(r,t)$ by
$(A,B)$; otherwise replace $\ell$ by $\ell-1$ and $(r,t)$ by $(B,A)$.
The new residual is positive, has both coordinates below $S$, and remains
in the cone. This last adjustment is essential: the translation theorem
was not asserted at the empty state $(0,0)$.
\end{proof}

\begin{theorem}[Tape optimality from finite translation and cone checks]
\label{thm:finite-cone-optimality}
Suppose $1\le A<B$, the hypotheses of Theorem~\ref{thm:translation} hold,
and the following finite equalities have been verified:
\begin{equation}
 D_{00}(r,t)=r+t-1
 \quad\text{for }1\le r\le t<S,\quad A t\le B r.
 \tag{C}
\end{equation}
Then
\[
 M(m,n)=m+n-1\qquad\text{whenever }m\le n\le\frac BA m.
\]
\end{theorem}
\begin{proof}
The unconstrained translation in Theorem~\ref{thm:translation} gives
\[
 D_{00}(u+A,v+B)\ge D_{00}(u,v)+S
\]
for every $u,v\ge0$ with $u+v\ge1$. Symmetry under swapping the two lists
also gives
\[
 D_{00}(u+B,v+A)\ge D_{00}(u,v)+S.
\]
For the desired $(m,n)$, both inequalities defining the cone hold.
Use Lemma~\ref{lem:cone-remainder}. By (C) and list-swap symmetry,
$D_{00}(r,t)=r+t-1$ also when $r>t$. Starting with this positive residual,
iterate the first translation $k$ times and the second $\ell$ times.
All intermediate dimensions are positive. It follows that
\[
 D_{00}(m,n)\ge(k+\ell)S+r+t-1=m+n-1.
\]
The adversary lower bound $D_{00}\le M$, established by the game-semantics
argument, and the ordinary tape-merge upper bound $M\le m+n-1$ force equality.
No diagonal-extension, constrained-increment, or deficit-monotonicity
lemma is used.
\end{proof}

\subsection{The concrete parameters and the finite obligations}
For $(A,B,H)=(95,148,152)$, we have $S=243$ and
$H\ge\max\{A+4,B+3\}$. The finite square uses indices no larger than
$(247,300)$ and so is contained in a table computed through dimensions 300.
The seven boundary thresholds in (B), in their displayed order, are
\[
 242,\quad242,\quad243,\quad242,\quad243,\quad244,\quad244.
\]
The cone check (C) consists of 10,651 pairs. Equivalently, checking both
orders uses 21,060 pairs, since the 242 diagonal pairs are counted once
rather than twice. The finite translation square has 138,622 feasible,
nonexcluded inequalities. These exact finite obligations, together with
the recurrence definition and the induction just proved, imply
\[
 \boxed{M(m,n)=m+n-1\quad\text{for}\quad m\le n\le\frac{148}{95}m.}
\]
The separate computation section records how each obligation is evaluated
and independently checked. 


\section{The finite computation and reproduction}
\label{sec:computation}

\subsection{Evaluating the recurrence exactly}
We evaluated all nine states with $1\le m,n\le300$, in increasing order
of $m+n$. Every child in a legal response has smaller total size, so each
lookup refers to a value already proved by the recurrence. There are
$9\cdot300^2=810\,000$ stored entries, including 600 impossible states.
All arithmetic is integral. Impossible states are represented in the files
by $-10000$; a response with such a child is discarded before addition.
Every feasible value lies between 1 and $m+n-1\le599$ by
\eqref{eq:disjunctive-soundness}, so the signed 16-bit storage of the primary
program is sufficient. Zero-axis values are supplied by the explicit
normalization rules of Section~\ref{sec:game-definition}.

The primary program, \texttt{adversary/disjunctive\_fast.cpp}, uses the
three geometric split forms. Each split supplies its child sum to one or
two rectangles of possible first queries. For example, a simple cut
$(p,q)$ applies to $i\le p,j\ge q+1$ and $i\ge p+1,j\le q$.
Store the sum at the corner of each applicable rectangle, then propagate
maxima in the corresponding two coordinate directions. The resulting
entry at $(i,j)$ is exactly the greatest response sum available there.
Taking one plus the least of these entries implements
\eqref{eq:nine-state-recurrence}. Enumerating split corners and propagating
rectangle maxima both cost $O(mn)$ operations per boundary state, rather
than enumerating every split separately for every query. Summing over
dimensions at most $N$ gives $O(N^4)$ operations and $O(N^2)$ stored values;
the factors for the nine boundary types are constant.

A separate implementation, \texttt{dp\_audit/fast\_literal.cpp}, was
written from the six published formulas without consulting the primary
source. It computes all nine boundary types without symmetry reduction.
It keeps six rectangle-max arrays, one per formula. In the order of the
rows in Section~\ref{sec:game-definition}, their query entries are
$(i,j),(i,j-1),(i+1,j),(i,j),(i,j+1),(i-1,j)$.
The strict offsets enforce the pivot exclusions. Cells that would refer
to the current parent but cannot be queried are omitted.
The independently generated complete tables agree byte for byte.
A slower literal implementation, \texttt{dp\_audit/literal.cpp}, enumerates
all six parameter ranges directly. Its 15,912 entries on
$1\le m\le34,1\le n\le52$ also agree with the primary table.

The common complete CSV file has the SHA-256 digest
\begin{center}\small\ttfamily
254a1a1c9efc1b7efe96585934ba7138\\
57840f4230bc9388874017845739307a
\end{center}
The two displayed lines concatenate to one digest. A digest identifies the
data; it is not a substitute for computing or checking the recurrence.
The CSV columns are \texttt{m,n,l,r,value}, with $0,L,R$ encoded as $0,1,2$.

\subsection{The checked hypotheses}
The final checker, \texttt{audit\_arxiv/verify\_main\_theorem.py}, first
validates the table's dimensions, distinct row keys, feasibility sentinels,
and integer ranges. It then checks exactly (B), (T), and (C), with
$(A,B,H)=(95,148,152)$.
The seven boundary values, all attaining their required thresholds, are
\begin{center}
\begin{tabular}{lrr}
\toprule
State and dimensions & Computed value & Required minimum\\
\midrule
$D_{00}(95,148)$ &242&242\\
$D_{R0}(95,148)$ &242&242\\
$D_{R0}(95,149)$ &243&243\\
$D_{LR}(96,148)$ &242&242\\
$D_{L0}(96,148)$ &243&243\\
$D_{L0}(97,148)$ &244&244\\
$D_{L0}(96,149)$ &244&244\\
\bottomrule
\end{tabular}
\end{center}
For each translation define its slack to be the left side of (T) minus
the right side. The counts below include only the stated feasible,
nonexcluded domains.
\begin{center}
\begin{tabular}{lrr}
\toprule
Boundary type & Inequalities checked & Minimum slack\\
\midrule
$00$ &23,408&0\\
$L0$ &23,104&0\\
$R0$ &23,104&0\\
$RL$ &22,952&0\\
$LL$ &23,103&0\\
$LR$ &22,951&0\\
\midrule
Total&138,622&0\\
\bottomrule
\end{tabular}
\end{center}
The excluded translations from $LL(1,1)$ and $LR(2,1)$ each have slack
$-1$. Thus omitting their repairs would invalidate the proof.
Finally the checker tests all 10,651 pairs
$1\le r\le t\le242$ satisfying $95t\le148r$ and finds
$D_{00}(r,t)=r+t-1$ in every case. List-swap symmetry covers the reverse
ordering of the residual coordinates.
Theorem~\ref{thm:finite-cone-optimality} now proves
Theorem~\ref{thm:main}.

\subsection{What is and is not certified}
The finite program proves integer statements about the recurrence.
Section~\ref{sec:game-definition} proves that the recurrence supplies a
lower bound for ordinary merging; Section~\ref{sec:expanded-translation}
proves the passage from finite checks to all dimensions. None of these
three steps can replace either of the others. In particular a value
$D_{00}(m,n)<m+n-1$ does not prove that tape merge is suboptimal: the
restricted adversary may simply be too weak.

The accompanying code is a reproducible computational certificate, not a
proof formalized in a proof assistant. The independent implementations,
literal cross-checks, and separate certificate verifier reduce the risk of
an implementation error. Readers can rerun the finite checks or regenerate
the tables from the displayed recurrences. No floating-point estimates or
numerically inferred asymptotic regularities enter the proof.


% Mathematical fragment for inclusion in the main manuscript.
% Requires amsmath, amssymb and amsthm. All labels have the prefix ex:.
\section{An unrestricted finite certificate for exact merging}
\label{ex:section}

We work with deterministic comparison decision trees and distinct keys.
The input consists of two internally ordered chains
$a_1<\cdots<a_m$ and $b_1<\cdots<b_n$.
A partial order $P$ records the two chain orders and all relations implied
by the comparisons already answered. A linear extension of $P$ is a total
order consistent with all these relations. Two keys are incomparable if
neither order between them is implied. At an internal node the algorithm may compare any two keys; comparisons whose
outcomes are already implied can be deleted.  Thus it suffices to consider
incomparable pairs, necessarily with one key in each chain.  The cost is the
maximum root-to-leaf number of comparisons.  Let $H(P)$ denote the optimum
remaining cost for a partial order $P$ obtained from the two chains and an
arbitrary consistent comparison history.  In particular,
$M(m,n)=H(P_{m,n})$ for the disjoint union of the two chains.
All lower-bound arguments below concern this unrestricted model.

\subsection{Exact states and their legal queries}

For a total order extending the two chains, let $x_i$ be the number of b-keys
preceding $a_i$.  The total order is uniquely specified by
\[
             0\le x_1\le\cdots\le x_m\le n.
\]
A state is a pair of nondecreasing integer vectors
$\ell=(\ell_1,\ldots,\ell_m)$ and $u=(u_1,\ldots,u_m)$ satisfying
\begin{equation}\label{ex:valid}
 0\le\ell_i\le u_i\le n\qquad(1\le i\le m).
\end{equation}
It represents exactly the rank vectors with $\ell_i\le x_i\le u_i$.
Equivalently, it contains the cross-chain relations
\[
 b_j<a_i\quad(j\le\ell_i),\qquad
 a_i<b_j\quad(j>u_i).
\]
The original merging problem has $\ell_i=0$ and $u_i=n$.
For example, with $m=n=2$, the rank vector $(0,2)$ describes
the order $a_1<b_1<b_2<a_2$. The answer $a_1<b_2$ restricts
$x_1$ to at most 1, while $b_2<a_1$ restricts it to at least 2.

\begin{lemma}[State sufficiency and query completeness]\label{ex:state-lemma}
Every consistent history of comparisons between the two chains has the above
representation.  In a represented state, $a_i$ and $b_j$ are incomparable
if and only if $\ell_i<j\le u_i$.  Both outcomes of each such comparison are
feasible.  Its two successor states are
\begin{align}
 u_k^-&=\begin{cases}\min(u_k,j-1),&k\le i,\\u_k,&k>i,\end{cases}
 &\ell^-&=\ell,\label{ex:minus}\\
 \ell_k^+&=\begin{cases}\ell_k,&k<i,\\\max(\ell_k,j),&k\ge i,\end{cases}
 &u^+&=u.\label{ex:plus}
\end{align}
The minus successor means $a_i<b_j$ and the plus successor means $b_j<a_i$.

\end{lemma}
\begin{proof}
The outcomes of a comparison impose $x_i\le j-1$ or $x_i\ge j$.
Monotonicity propagates these inequalities precisely as in
\eqref{ex:minus}--\eqref{ex:plus}; it introduces no disjunctive restriction.
Induction from the unrestricted initial state proves sufficiency.
Conversely, each feasible rank vector constructs one consistent total order,
by putting $a_i$ in the gap after $b_{x_i}$, retaining the order of the a-keys
when several ranks agree.  The vectors $\ell$ and $u$ themselves are feasible.
For an incomparable query, $\ell_i<j$ and $u_i\ge j$ show that these two
vectors witness the two possible outcomes.  The updated endpoints remain
nondecreasing and satisfy \eqref{ex:valid}; for example, when $k\le i$,
$\ell_k\le\ell_i\le j-1$ in the minus case.
All comparisons outside the stated range have fixed outcomes, as do all
within-chain comparisons, so the listed queries are exhaustive.
\end{proof}

Write
\begin{equation}\label{ex:area}
             \eta(P)=\sum_{i=1}^m(u_i-\ell_i).
\end{equation}
This is exactly the number of incomparable pairs.  Each legal query strictly
decreases $\eta$ in both successors: at least the queried pair becomes comparable,
and no known relation is removed.

Swapping the chains and reversing the whole total order preserve $H$.
For reversal the endpoints are
$\ell'_i=n-u_{m+1-i}$ and $u'_i=n-\ell_{m+1-i}$.
After swapping the chains, the interval for $b_j$, now in the first chain,
has endpoints
\begin{equation}\label{ex:transpose}
 \ell'_j=\#\{i:u_i<j\},\qquad u'_j=\#\{i:\ell_i<j\}.
\end{equation}
Here $\#$ means the number of indices satisfying the stated condition.
These operations commute and generate four equivalent representations.
The implementation takes their lexicographically least encoding (the first in dictionary order); this is
only a symmetry reduction and imposes no restriction on the query set.

\subsection{Ordinal-sum decomposition}

Let $G(P)$ be the graph whose vertices are the keys and whose edges are the
incomparable pairs. A connected component is a maximal set of vertices joined
by paths of edges; it is nontrivial if it contains at least two vertices.  The neighbors of $a_i$ are exactly
$b_{\ell_i+1},\ldots,b_{u_i}$.  The nontrivial connected components can therefore
be found by grouping successive nonempty row intervals whenever they overlap.
A group beginning at row $r$ and ending at row $s$ has second-chain indices
$\ell_r+1,\ldots,u_s$; its normalized endpoints are
$\ell_i-\ell_r,u_i-\ell_r$, for $r\le i\le s$.
A new group begins if the next row is empty or its left endpoint is at least
the preceding right endpoint.  Monotonicity of both endpoint sequences
ensures that no later row can bridge such a separation.  Isolated vertices
are retained in the determined output order and contribute zero cost.

\begin{lemma}[Component additivity]\label{ex:component-lemma}
If $P_1,\ldots,P_t$ are the nontrivial components just described, then
\begin{equation}\label{ex:additivity}
                    H(P)=\sum_{r=1}^t H(P_r).
\end{equation}

\end{lemma}
\begin{proof}
First, connected components of an incomparability graph are ordinal-sum
factors.  Indeed, vertices from different components are comparable.  For a
fixed vertex $y$ outside a component, its comparison direction with vertices
of that component cannot change along an incomparability edge: $x<y<z$
would imply $x<z$.  Propagating along a path proves that $y$ lies entirely
before or entirely after the component.  Applying this to both components
shows that their direction is uniform.  These directions give a linear order
of the components.  Their internal linear extensions can consequently be
chosen independently.

This ordered concatenation of independently ordered parts is called an
ordinal sum. An upper bound is obtained by solving the components separately.  For the
reverse inequality, queries between different components have known answers
and can be discarded.  In each component, the adversary can maintain its
local optimal remaining height.  For any query in a component of height $h$,
some answer leaves local height at least $h-1$; otherwise that query would
give an algorithm of height less than $h$.  The adversary uses this answer in
the queried component and leaves all others unchanged.  The sum of the local
heights decreases by at most one per comparison and must be zero when the
whole order is determined.  Initially it is the right side of
\eqref{ex:additivity}.  This proves the lower bound even for algorithms that
interleave queries among components.
\end{proof}

There is no nontrivial component exactly when $P$ is a total order, in which
case $H(P)=0$.  The sum of the areas of all components is $\eta(P)$.  In
particular, every component of either successor of a legal query has area
strictly smaller than that of the original state.

\subsection{Elementary bounds with complete recurrences}

Let $e(P)$ be the number of linear extensions. The indicator $\mathbf 1_E$ is 1 when the condition $E$ holds and 0 otherwise.  It is computed without
enumerating them by setting $d_0(0)=1$, $d_0(j)=0$ for $j>0$, and
\begin{equation}\label{ex:count}
 d_i(j)=\mathbf 1_{\{\ell_i\le j\le u_i\}}
                  \sum_{r=0}^j d_{i-1}(r),
 \qquad e(P)=\sum_{j=0}^n d_m(j).
\end{equation}
Here $d_i(j)$ counts feasible rank prefixes ending at $x_i=j$; the recurrence
follows by conditioning on $x_{i-1}$.  Thus the decision-tree leaf bound gives
\begin{equation}\label{ex:information}
                     H(P)\ge I(P):=\lceil\log_2 e(P)\rceil.
\end{equation}

For a linear extension $\pi$ of $P$, let $w_P(\pi)$ be the number of adjacent
pairs in $\pi$ that are incomparable in $P$.  Such a pair necessarily consists
of keys from different chains.  Set $W(P)=\max_\pi w_P(\pi)$.

\begin{lemma}[Unresolved adjacency bound]\label{ex:adjacency-lemma}
For every represented state $P$, one has $H(P)\ge W(P)$.

\end{lemma}
\begin{proof}
Fix an extension $\pi$.  If adjacent keys $x,y$ in $\pi$ are incomparable in
$P$, swapping them gives another extension of $P$.  Every comparison other
than the comparison of $x$ and $y$ has the same outcome on these two inputs.
A correct algorithm must therefore compare $x$ with $y$ on its execution on
$\pi$.  This applies separately to every counted adjacent pair, and proves
that this execution has length at least $w_P(\pi)$.  Maximizing over $\pi$
proves the lemma.  Pairs already comparable in $P$ are not counted.
\end{proof}

Here are explicit polynomial-time recurrences for $W$ and for an elementary
upper bound.  A prefix of a merged order is described by the numbers $(i,j)$
of emitted a- and b-keys.  Put $\ell_0=0$ and $u_{m+1}=n$.
A feasible prefix vertex satisfies $\ell_i\le j\le u_{i+1}$.
From a feasible vertex, the permissible next moves are
\begin{align*}
 \mathcal A(i,j)&:\quad i<m,\quad \ell_{i+1}\le j\le u_{i+1},\\
 \mathcal B(i,j)&:\quad j<n,\quad \ell_i\le j+1\le u_{i+1}.
\end{align*}
Every root-to-terminal path using these moves corresponds to one extension
of $P$, and every extension gives such a path.  This follows either by the
rank-vector bijection or by checking the required rank of each newly emitted
a-key.  Only vertices reachable from $(0,0)$ are needed.

Let $F_A(i,j)$ and $F_B(i,j)$ be the maximum future contribution to
$w_P(\pi)$, given a prefix ending in an a-key or a b-key, respectively.
The switch weights are
\[
 \alpha(i,j)=\mathbf 1_{\{j>0,\ \ell_{i+1}<j\le u_{i+1}\}},\qquad
 \beta(i,j)=\mathbf 1_{\{i>0,\ \ell_i<j+1\le u_i\}}.
\]
The first is used only for a permissible $\mathcal A$ move, and the second
only for a permissible $\mathcal B$ move.  With both terminal values zero,
backward dynamic programming gives
\begin{align}
 F_A(i,j)&=\max\left(
 \{F_A(i+1,j):\mathcal A(i,j)\},
 \{\beta(i,j)+F_B(i,j+1):\mathcal B(i,j)\}\right),\label{ex:wa}\\
 F_B(i,j)&=\max\left(
 \{\alpha(i,j)+F_A(i+1,j):\mathcal A(i,j)\},
 \{F_B(i,j+1):\mathcal B(i,j)\}\right).\label{ex:wb}
\end{align}
An unavailable option is omitted from the maximum.  At the empty prefix the
switch weights for the first move are zero, and hence
\[
                     W(P)=\max\{F_A(0,0),F_B(0,0)\}.
\]
These formulas explicitly count only unresolved cross-chain adjacencies.

For an upper bound, use ordinary front-element merging, omitting a comparison
when the next key is already forced.  If exactly one of the two moves is
permissible, emit that next key for free.  If both are permissible, compare
the two front keys and emit the smaller one.  Let $T(m,n)=0$, and set
\begin{equation}\label{ex:tape}
 T(i,j)=\mathbf 1_{\{\mathcal A(i,j),\mathcal B(i,j)\}}
       +\max\left(\{T(i+1,j):\mathcal A(i,j)\},
                   \{T(i,j+1):\mathcal B(i,j)\}\right)
\end{equation}
for a nonterminal feasible prefix.  Then $H(P)\le T(0,0)$.
The comparison just used involves the emitted key, so it introduces no
additional restriction among the un-emitted keys beyond the original state
and the enlarged prefix.  Consequently \eqref{ex:tape} is the exact
worst-case cost of this specified algorithm.

For a nontrivial indecomposable state with one chain of size one, every key
of the other chain is incomparable with it.  There are exactly $n+1$
possible insertion gaps.  Binary insertion and the leaf bound give the exact
value $\lceil\log_2(n+1)\rceil$, with the analogous formula after exchanging
the chains.  We use this exact value in place of the front-merge upper bound.
Thus the elementary bounds $E_-(P),E_+(P)$ used below are
\[
 E_-(P)=\max\{I(P),W(P)\},\qquad E_+(P)=T(0,0),
\]
with this one-element-chain refinement.  For a decomposable state, we apply
the bounds to its nontrivial components and sum them.

\subsection{A finite, well-founded certificate}

For a nonterminal state, the unrestricted decision-tree recurrence is
\begin{equation}\label{ex:bellman}
 H(P)=1+\min_{\ell_i<j\le u_i}
               \max\{H(P^-_{ij}),H(P^+_{ij})\}.
\end{equation}
The subscripts $ij$ denote the successors of query $a_i:b_j$. The root of any decision tree makes one of the listed queries, and its two
subtrees solve the two successors.  Conversely, combining a query with
optimal successor trees constructs a valid tree.  This proves
\eqref{ex:bellman}, including both inequalities.

A certificate assigns integer bounds $L(P),U(P)$ to finitely many canonical,
nontrivial indecomposable states.  All other such states receive the
explicit defaults $E_-(P),E_+(P)$.  For any state $Q$, define
\[
 \mathcal L(Q)=\sum_{C\in\operatorname{comp}(Q)}L(C),\qquad
 \mathcal U(Q)=\sum_{C\in\operatorname{comp}(Q)}U(C),
\]
where $\operatorname{comp}(Q)$ denotes the set of nontrivial normalized
components, and a missing component record is interpreted using the defaults.
The empty sum is zero.  Each explicit record must be a valid canonical state
and satisfy $0\le L(P)\le U(P)$.  Its two assertions are checked as follows:
\begin{enumerate}
 \item If $L(P)>E_-(P)$, check \emph{every} legal query and require
 \begin{equation}\label{ex:checklower}
  1+\max\{\mathcal L(P^-_{ij}),\mathcal L(P^+_{ij})\}\ge L(P).
 \end{equation}
 \item If $U(P)<E_+(P)$, check that \emph{at least one} legal query satisfies
 \begin{equation}\label{ex:checkupper}
  1+\max\{\mathcal U(P^-_{ij}),\mathcal U(P^+_{ij})\}\le U(P).
 \end{equation}
\end{enumerate}
When a proposed bound is weaker than its corresponding elementary bound,
no Bellman check is needed for that bound.

\begin{theorem}[Certificate soundness]\label{ex:certificate-theorem}
If these checks pass, every explicit and default assertion is true:
$L(P)\le H(P)\le U(P)$.

\end{theorem}
\begin{proof}
Induct on the nonnegative integer $\eta(P)$ in \eqref{ex:area}.
Defaults are true by the elementary bounds, regardless of the area.
For an explicit assertion not established by a default bound, every component
of every successor appearing in \eqref{ex:checklower} or
\eqref{ex:checkupper} has strictly smaller area.  By the induction hypothesis
and Lemma~\ref{ex:component-lemma} its aggregated bounds are true.  Equation~\eqref{ex:bellman} now
proves the requested lower or upper bound.  The base case is a total order,
with area and cost zero.  This also shows that it is unnecessary to list or
verify the certificate records in topological order: their logical dependency
is nevertheless acyclic.  Missing records cannot create an unsupported
assertion because their defaults have independent proofs.
\end{proof}

The upper certificate is constructive.  At an explicit state with an improved
upper bound, enumerate its queries and use a witness to
\eqref{ex:checkupper}.  After the observed outcome, solve its ordinal-sum
components recursively.  At a default upper bound use the front algorithm
\eqref{ex:tape}, or binary insertion in the one-element case.  Canonicalization
only requires tracking the chain exchange and order reversal when translating
indices and outcomes back to the current keys.  Strict area decrease proves
termination, and the checked inequalities prove the asserted worst-case cost.
The lower certificate similarly records, implicitly or explicitly, a hard
outcome for every legal query.  Its elementary leaves are justified by
Lemma~\ref{ex:adjacency-lemma} or the decision-tree leaf count, not by a restriction on allowed
algorithms.

\subsection{The two checked instances}

Applying this procedure to the supplied integer certificates yields
\begin{equation}\label{ex:results}
                         M(6,22)=21,
                 \qquad M(6,25)=22.
\end{equation}
The finite verification statistics are
\[
\begin{array}{c|r|r|c}
(m,n)&\text{explicit records}&\text{default states evaluated}
    &\text{certified root interval}\\\hline
(6,22)&454\,963&5\,560\,670&[21,21]\\
(6,25)&1\,111\,896&13\,522\,983&[22,22].
\end{array}
\]
The verifier is a standalone program: it neither imports nor runs the search
program.  It evaluates \eqref{ex:count}, the backward recurrences
\eqref{ex:wa}--\eqref{ex:tape}, and all the required Bellman inequalities.
The readable verifier validates endpoint inequalities, monotonicity,
dimensions, indecomposability, integer ranges, canonical form and uniqueness
of every explicit record before checking its assertions.  Its requirements
use exceptions and cannot be disabled by compiler assertion flags.  Both full
passes were also run with assertions disabled.  A separate streaming check
independently validates the endpoint and range contracts for both files.  All encountered states have at most 31 keys and at most
six keys in the shorter chain.  Their extension counts are at most
$\binom{31}{6}=736\,281$, so the exact unsigned 64-bit arithmetic used for
counting cannot overflow in these certificates.  The largest asserted height
is 30, and all endpoint and height encodings fit their specified byte fields.
No published numerical merging values are used as default bounds or seeds.

As an additional software check, an independently written solver on explicit
sets of interleavings, using neither component decomposition nor symmetry
reduction nor the adjacency bound, agrees on all 940 consistent interval
states with $1\le m\le3$ and $1\le n\le4$.  This finite cross-check supplements,
and does not replace, the proofs of the reductions and the complete
certificate verification.

For $(6,22)$, a certified optimal first query is $a_3$ versus $b_{10}$.
The two constrained successors both have upper bound 20; their reflected
counterpart gives the query $a_4$ versus $b_{13}$.  The lower certificate
checks all $6\cdot22=132$ possible first queries and in each case certifies
an outcome requiring at least 20 more comparisons.  These are certified
optimal first queries, without a claim that they exhaust all optimal choices.


\section{Further directions}
The main result uses a larger block within the existing disjunctive
framework. It therefore leaves two different questions open: how far
block translations can be extended within that framework, and how much
stronger a different adversary can be. The obstruction proved by
Li--Sun--Zhang near ratio $1.8$ concerns this particular lower-bound
method, not an algorithm proving tape merge suboptimal at that ratio.

Larger-looking numerical blocks must be checked against the complete
translation hypothesis, including small constrained states. For example,
the candidate blocks $(91,142)$ and $(93,145)$ both fail the proposed
translation at $RL(1,2)$, $LL(2,2)$, and $LR(3,2)$, each with slack $-1$.
Thus an improved ratio cannot be inferred from a few favorable block
boundary values. Additional exceptional-state repairs, a different
translation scheme, or a stronger game would be needed.

The exact certificates illustrate a complementary direction. Their lower
bounds quantify over every legal comparison, while their upper bounds
describe actual algorithms. They settle the two finite equalities in
Theorem~\ref{thm:exact}, but do not establish a formula for all $M(6,n)$.
In particular, Smith--Lang's further implications involving $(8,29)$ also
need a separate assertion about $(7,22)$; the two equalities proved here
alone do not resolve that case. Neither the exact computations nor the
upper-bound appendix solve the global merging problem.

\appendix


% Appendix fragment. Requires amsmath, amsthm and standard lemma/proposition
% environments. Bibliography keys: smithlang, lsz2016.
\section{An upper-bound consequence of the modified binary merge}
\label{app:upper-consequence}

This appendix records a consequence of the finite algorithms of
Smith--Lang~\cite{smithlang} and the modified binary merge of
Li--Sun--Zhang~\cite{lsz2016}.  Smith--Lang, Section~5.1, already
conjectured the equalities
\[
 M(m,2m-3+k)=3m-5+k\qquad(m\ge9,\ 0\le k\le2).
\]
We prove their upper-bound direction.  We make no claim here about the
matching lower bounds or the priority of this consequence.  The recursive
argument is included in full; the finite comparison algorithms explicitly
listed below are used as established external results.

We put $M(m,0)=M(0,n)=0$.  All list elements are distinct, and all comparisons
are comparisons between elements of the two initially sorted lists.

\begin{lemma}[Elementary bounds]\label{lem:upper-elementary}
For $m,n\ge1$,
\[
 M(m,n)\le m+n-1.
\]
For $m\ge1$ and $n\ge2$,
\begin{equation}\label{eq:upper-pair-bound}
 M(m,n)\le 2m+\lfloor n/2\rfloor-1.
\end{equation}
\end{lemma}
\begin{proof}
The first assertion is tape merge.  For the second, first merge the $m$
a-elements with the $\lfloor n/2\rfloor$ even-indexed b-elements
by tape merge.  This uses at most $m+\lfloor n/2\rfloor-1$ comparisons.
Each gap between consecutive even-indexed elements, including the two end
gaps, contains at most one omitted odd-indexed b-element.  If that gap
contains $t$ a-elements, insert its omitted element using at most $t$
comparisons.  Summing over the gaps costs at most $m$ more comparisons.
\end{proof}

\begin{lemma}[Li--Sun--Zhang recurrence]\label{lem:upper-lsz-recurrence}
For $m,n\ge2$,
\begin{equation}\label{eq:upper-lsz-recurrence}
\begin{split}
 M(m,n)\le\max\{&M(m,n-2)+1,\ M(m-1,n-2)+3,\\
                 &M(m-2,n)+3,\ M(m-2,n-1)+4\}.
\end{split}
\end{equation}
\end{lemma}
\begin{proof}
Compare $a_1$ with $b_2$.  If $b_2<a_1$, output $b_1,b_2$ and continue with
sizes $(m,n-2)$, having spent one comparison.
Otherwise compare $a_2$ with $b_2$.  If $b_2<a_2$, compare $a_1$ with $b_1$;
the three-element prefix consisting of $a_1,b_1,b_2$ is then ordered, leaving
sizes $(m-1,n-2)$ after three comparisons.
In the remaining case $a_2<b_2$, compare $a_2$ with $b_1$.  If $a_2<b_1$,
output $a_1,a_2$, leaving $(m-2,n)$ after three comparisons.  If $b_1<a_2$,
compare $a_1$ with $b_1$ and output the ordered prefix consisting of
$a_1,a_2,b_1$, leaving $(m-2,n-1)$ after four comparisons.
Any additional information about the remaining lists can be discarded.
\end{proof}

\paragraph{Finite inputs.}
The following twelve finite upper bounds are recorded in the table of
Smith--Lang~\cite{smithlang}; several had been established earlier.
Only upper bounds are needed.
\begin{equation}\label{eq:upper-finite-seeds}
\begin{gathered}
\begin{array}{c|rrrr}
 (m,n)&(3,6)&(3,8)&(3,10)&(3,12)\\ \hline
 \text{upper bound}&7&8&9&10
\end{array}
\\[6pt]
\begin{array}{c|rrr}
 (m,n)&(5,9)&(5,11)&(5,13)\\ \hline
\text{upper bound}&12&13&14
\end{array}
\\[6pt]
\begin{array}{c|rrrrr}
 (m,n)&(7,12)&(10,20)&(9,15)&(10,17)&(11,19)\\ \hline
\text{upper bound}&17&27&22&25&28.
\end{array}
\end{gathered}
\end{equation}
We do not regard the published abbreviated drawings, whose leaves can refer
to omitted subtrees, as complete independently verified decision-tree
certificates.  Thus the finite assertions in~\eqref{eq:upper-finite-seeds}
are the precise external input to this appendix.

\begin{lemma}[The finite strips needed below]\label{lem:upper-strips}
The bounds
\begin{align}
 M(m,2m+2k)&\le3m+k-2
   &&(m\ge3,\quad k\in\{0,1,2,3\}),\label{eq:upper-even-strip}\\
 M(m,2m+2k-1)&\le3m+k-3
   &&(m\ge5,\quad k\in\{0,1,2\})\label{eq:upper-odd-strip}
\end{align}
follow from the finite inputs and Lemmas~\ref{lem:upper-elementary}
and~\ref{lem:upper-lsz-recurrence}.
\end{lemma}
\begin{proof}
For~\eqref{eq:upper-even-strip}, use the four $m=3$ inputs and induction on
$m+k$ for $m\ge4$.  Apply~\eqref{eq:upper-lsz-recurrence} at $n=2m+2k$.
Its first term is at most $3m+k-2$ by the induction hypothesis with $k-1$
when $k\ge1$.  When $k=0$, tape merge gives that same bound:
$M(m,2m-2)+1\le3m-2$.  The second term is at most
$3(m-1)+k-2+3=3m+k-2$ by induction on $m$.
For the last two terms, put $r=m-2\ge2$ and use
\eqref{eq:upper-pair-bound}:
\begin{align*}
 M(r,2r+4+2k)+3
   &\le 2r+(r+2+k)-1+3=3m+k-2,\\
 M(r,2r+3+2k)+4
   &\le 2r+(r+1+k)-1+4=3m+k-2.
\end{align*}
This proves the even strip in its stated finite range of $k$.

For~\eqref{eq:upper-odd-strip}, use the three $m=5$ inputs and again induct
on $m+k$, now for $m\ge6$.  In the first term of the recurrence use the
odd-strip induction hypothesis with $k-1$ if $k\ge1$; at $k=0$ use
$M(m,2m-3)+1\le3m-3$ from tape merge.  The second term uses the odd-strip
hypothesis at $m-1$.  Both are at most $3m+k-3$.
Set $r=m-2\ge4$ in the remaining terms.  The elementary pair bound gives
\[
 M(r,2r+3+2k)+3\le 3r+k+3=3m+k-3.
\]
The already proved even strip, with index $k+1\in\{1,2,3\}$, gives
\[
 M(r,2r+2+2k)+4\le3r+(k+1)-2+4=3m+k-3.
\]
All dimensions are positive and all applications lie in the displayed
parameter ranges.
\end{proof}

\begin{lemma}[Two further diagonals]\label{lem:upper-two-diagonals}
One has
\begin{align}
 M(m,2m-2)&\le3m-4 &&(m\ge7),\label{eq:upper-minus-two}\\
 M(m,2m)&\le3m-3 &&(m\ge10).\label{eq:upper-zero}
\end{align}
\end{lemma}
\begin{proof}
The first inequality starts from $M(7,12)\le17$.
For $m\ge8$, apply the recurrence with $n=2m-2$.
Tape merge, induction, the even strip with index $1$, and the odd strip
with index $1$, respectively, bound its four terms by
\[
\begin{split}
 M(m,2m-4)+1&\le3m-4,\\
 M(m-1,2m-4)+3&\le3(m-1)-4+3=3m-4,\\
 M(m-2,2m-2)+3&\le3(m-2)+1-2+3=3m-4,\\
 M(m-2,2m-3)+4&\le3(m-2)+1-3+4=3m-4.
\end{split}
\]
The strip applications have $m-2\ge6$, so their size conditions hold.

The second inequality starts from $M(10,20)\le27$.
For $m\ge11$, apply the recurrence with $n=2m$.
Use~\eqref{eq:upper-minus-two}, induction, the even strip with index $2$,
and the odd strip with index $2$:
\[
\begin{split}
 M(m,2m-2)+1&\le3m-3,\\
 M(m-1,2m-2)+3&\le3(m-1)-3+3=3m-3,\\
 M(m-2,2m)+3&\le3(m-2)+2-2+3=3m-3,\\
 M(m-2,2m-1)+4&\le3(m-2)+2-3+4=3m-3.
\end{split}
\]
Here $m-2\ge9$, again within the strip domains.
\end{proof}

Define
\[
 \alpha(m)=\max\{n\ge m:M(m,n)=m+n-1\}.
\]
The set is nonempty since tape merge is optimal at $n=m$, and finite
because \eqref{eq:upper-pair-bound} is strictly below tape cost for $n\ge2m+1$.

\begin{proposition}[Upper direction of the Smith--Lang pattern]
\label{prop:upper-smith-lang}
For every integer $m\ge9$ and $k\in\{0,1,2\}$,
\begin{equation}\label{eq:upper-smith-lang}
 M(m,2m-3+k)\le3m-5+k.
\end{equation}
In particular, tape merge is not optimal whenever $m\ge9$ and $n\ge2m-3$.
Consequently $\alpha(m)\le2m-4$ for $m\ge9$.
\end{proposition}
\begin{proof}
For $k=0$, the cases $m=9,10,11$ are the final three finite inputs.
Suppose $m\ge12$ and use the recurrence at $n=2m-3$.
Tape merge, induction on $m$, the odd strip with index $1$, and
\eqref{eq:upper-zero}, respectively, give
\[
\begin{split}
 M(m,2m-5)+1&\le3m-5,\\
 M(m-1,2m-5)+3&\le3(m-1)-5+3=3m-5,\\
 M(m-2,2m-3)+3&\le3(m-2)+1-3+3=3m-5,\\
 M(m-2,2m-4)+4&\le3(m-2)-3+4=3m-5.
\end{split}
\]
The last application is valid because $m-2\ge10$.
This proves the $k=0$ case.  The case $k=1$ is
\eqref{eq:upper-minus-two}, and $k=2$ is the odd strip with index $0$.

Each of the three bounds saves one comparison over tape merge.
At $n=2m$, the even strip with index $0$ gives
$M(m,2m)\le3m-2$, saving one comparison for every $m\ge9$.
For every $n\ge2m+1$, the elementary pair bound gives
\[
 M(m,n)\le2m+\lfloor n/2\rfloor-1\le m+n-2,
\]
since $\lceil n/2\rceil\ge m+1$.  Thus every $n\ge2m-3$ is excluded
from the tape-optimal region.  The stated bound on $\alpha(m)$ follows
directly from its definition, without a monotonicity conjecture for $M$.
\end{proof}

\paragraph{Scope.}
This argument leaves the asymptotic upper boundary at ratio $2$.
It does not prove the equalities conjectured by Smith--Lang.
The recurrence is the published Li--Sun--Zhang algorithm, and the finite
inputs are published results.  The contribution of this appendix is the
explicit propagation and its carefully specified scope.


\section{Reproducing the finite proof}
\label{app:reproduction}
The arXiv source contains a single self-contained \texttt{main.tex} and an
\texttt{anc/} directory of ancillary source code, data, and certificates.
The document requires only standard \LaTeX{} packages and compiles with
\texttt{pdflatex main.tex}, run until references stabilize. No shell escape,
external bibliography processor, or network access is needed.
For the commands below, change into \texttt{anc/}. Python~3 and a compiler
supporting C++17 suffice. \texttt{README.txt} gives the complete file map,
and \texttt{SHA256SUMS} records the accepted files' digests.

The quickest main-theorem check uses the supplied table:
\begin{lstlisting}
python3 audit_arxiv/verify_main_theorem.py adversary/table_300.csv
\end{lstlisting}
It must report \texttt{PASS}, seven boundary checks, 138,622 translation
inequalities, and 10,651 cone seeds. This checks the finite hypotheses; to
recompute their table independently, run
\begin{lstlisting}
g++ -O3 -std=c++17 dp_audit/fast_literal.cpp -o fast_literal
./fast_literal 300 300 > independent_table.csv
cmp adversary/table_300.csv independent_table.csv
python3 audit_arxiv/verify_main_theorem.py independent_table.csv
\end{lstlisting}
The \texttt{cmp} command must succeed without output. The independent
full computation took about five minutes on the machine used for this
work; this timing is descriptive and is not a complexity guarantee.
The primary program and the smaller literal cross-check are also included.

The exact-instance proofs can be checked without rerunning their searches:
\begin{lstlisting}
gzip -dk exact/certificate_6_22_safe.txt.gz
gzip -dk exact/certificate_6_25_safe.txt.gz
python3 audit_arxiv/check_state_validity.py
g++ -O3 -std=c++17 exact/verify_certificate.cpp -o verify_exact
./verify_exact exact/certificate_6_22_safe.txt 6 22
./verify_exact exact/certificate_6_25_safe.txt 6 25
\end{lstlisting}
The first Python check verifies the dimensions and integer contracts of
every record. The final two runs must report \texttt{verified:true},
root values 21 and 22, and respectively 454,963 and 1,111,896 explicit
records. Those runs took approximately 67 and 185 seconds, respectively,
and cache several million elementary states in memory. The original
verifier and supplementary Python check use assertions: do not compile
with \texttt{-DNDEBUG}, run Python with \texttt{-O}, or set
\texttt{PYTHONOPTIMIZE}. The standalone readable verifier additionally
performs the record-contract checks itself and is described in the README.

The certificate file format is one record per line:
\[
 m\quad n\quad L\quad U\quad
 \ell_1\ \cdots\ \ell_m\quad u_1\ \cdots\ u_m.
\]
Here $L,U$ are asserted comparison bounds, while the endpoint vectors
$\ell,u$ specify the state in Section~\ref{ex:section}. Records are
canonical nontrivial components, with no duplicates. There are no
embedded published merging values: missing states use exactly the
elementary bounds proved in that section. The upper witness query is
recovered by enumerating legal queries; a separate decision tree with
one explicit leaf per full execution is unnecessary.

The two uncompressed certificate digests are, respectively,
\begin{center}\small\ttfamily
7fff5828213ebcb112ba11f3685b8680\\
9baa852d62ed73632893b39d4f924412\par\medskip
1e14335f5c4d77021fd0f52df3650e6c\\
db360e11dd9096015260ad40ebd8450d
\end{center}
Concatenate the two lines within each digest. The ancillary manifest
also identifies the compressed files, source programs, and finite
check reports. Historical exploratory outputs and third-party papers
are not part of the accepted arXiv proof data.


\begin{thebibliography}{99}

\bibitem{christen1978}
C. Christen.
Improving the bounds on optimal merging.
In \emph{19th Annual Symposium on Foundations of Computer Science},
pp.~259--266. IEEE, 1978.

\bibitem{christenstraight}
C. Christen.
\emph{On the optimality of the straight merging algorithm}.
Unpublished manuscript, 1978. Bibliographic information and attribution
as reported in~\cite{lsz2016,lsz2020}.

\bibitem{fernandez1993}
Wenceslas Fern\'andez de la Vega, Sampath Kannan, and Mikl\'os Santha.
Two probabilistic results on merging.
\emph{SIAM Journal on Computing} 22(2):261--271, 1993.
\href{https://doi.org/10.1137/0222019}{doi:10.1137/0222019}.

\bibitem{hwang1980}
Frank K. Hwang.
Optimal merging of 3 elements with $n$ elements.
\emph{SIAM Journal on Computing} 9(2):298--320, 1980.
\href{https://doi.org/10.1137/0209026}{doi:10.1137/0209026}.

\bibitem{hwanglin1971}
Frank K. Hwang and Shen Lin.
Optimal merging of 2 elements with $n$ elements.
\emph{Acta Informatica} 1:145--158, 1971.
\href{https://doi.org/10.1007/BF00289521}{doi:10.1007/BF00289521}.

\bibitem{hwanglin1972}
Frank K. Hwang and Shen Lin.
A simple algorithm for merging two disjoint linearly ordered sets.
\emph{SIAM Journal on Computing} 1(1):31--39, 1972.
\href{https://doi.org/10.1137/0201004}{doi:10.1137/0201004}.

\bibitem{knuth}
Donald E. Knuth.
\emph{The Art of Computer Programming, Volume 3: Sorting and Searching}.
Second edition. Addison--Wesley, Reading, Massachusetts, 1998.
Section~5.3.2, pp.~197--207.
Author's corrections: \url{https://cs.stanford.edu/~knuth/taocp.html}.

\bibitem{lsz2016}
Qian Li, Xiaoming Sun, and Jialin Zhang.
On the optimality of tape merge of two lists with similar size.
In \emph{27th International Symposium on Algorithms and Computation
(ISAAC 2016)}, Leibniz International Proceedings in Informatics 64,
article~51, pp.~51:1--51:17, 2016.
\href{https://doi.org/10.4230/LIPIcs.ISAAC.2016.51}{doi:10.4230/LIPIcs.ISAAC.2016.51}.
Preprint: \href{https://arxiv.org/abs/1610.03266}{arXiv:1610.03266}.

\bibitem{lsz2020}
Qian Li, Xiaoming Sun, and Jialin Zhang.
On the optimality of tape merge of two lists with similar size.
\emph{Algorithmica} 82(7):2107--2132, 2020.
\href{https://doi.org/10.1007/s00453-020-00690-x}{doi:10.1007/s00453-020-00690-x}.

\bibitem{linial1984}
Nathan Linial.
The information-theoretic bound is good for merging.
\emph{SIAM Journal on Computing} 13(4):795--801, 1984.
\href{https://doi.org/10.1137/0213049}{doi:10.1137/0213049}.

\bibitem{manacher1979}
Glenn K. Manacher.
Significant improvements to the Hwang--Lin merging algorithm.
\emph{Journal of the ACM} 26(3):434--440, 1979.
\href{https://doi.org/10.1145/322139.322144}{doi:10.1145/322139.322144}.

\bibitem{murphy1978}
Paul E. Murphy.
\emph{A problem in optimal merging}.
Technical Report DCS-TR-69, Department of Computer Science,
Rutgers University, 1978.
\href{https://doi.org/10.7282/t3-24m2-c806}{doi:10.7282/t3-24m2-c806}.

\bibitem{murphypaull1979}
Paul E. Murphy and Marvin C. Paull.
Minimum comparison merging of sets of approximately equal size.
\emph{Information and Control} 42(1):87--96, 1979.
\href{https://doi.org/10.1016/S0019-9958(79)90186-4}{doi:10.1016/S0019-9958(79)90186-4}.

\bibitem{monting1981}
J\"urgen Schulte M\"onting.
Merging of 4 or 5 elements with $n$ elements.
\emph{Theoretical Computer Science} 14(1):19--37, 1981.
\href{https://doi.org/10.1016/0304-3975(81)90003-7}{doi:10.1016/0304-3975(81)90003-7}.

\bibitem{smithlang}
Warren D. Smith and Kevin J. Lang.
\emph{Values of the merging function and algorithm design as a game}.
Unpublished technical report, NEC Research Institute, 1994, with
accompanying table and decision-tree drawings.
\url{https://rangevoting.org/WarrenSmithPages/homepage/works.html}
(author's publication listing, with links to the report and table).

\bibitem{stockmeyeryao1980}
Paul K. Stockmeyer and F. Frances Yao.
On the optimality of linear merge.
\emph{SIAM Journal on Computing} 9(1):85--90, 1980.
\href{https://doi.org/10.1137/0209006}{doi:10.1137/0209006}.

\bibitem{thanhbui1982}
Mai Thanh and T. D. Bui.
An improvement of the binary merge algorithm.
\emph{BIT} 22(4):454--462, 1982.

\end{thebibliography}
\end{document}
